\section{Introduction}
\label{sec:introduction}

Temporal information is usually discussed relative to an external time standard: a laboratory clock defines when a signal arrives, how quickly a probe evolves, or which phase a waveform carries. For a fully autonomous quantum description that reference cannot be supplied for free. A finite clock, signal, controller, and readout must instead be treated as parts of one physical system. This raises a basic resource question: \emph{what must a finite autonomous system possess in order to carry information about relative time at a given frequency?}

The question is not answered by global time-translation asymmetry alone. In Page--Wootters-type relational descriptions, nontrivial relative evolution can occur inside a state that is stationary under the total Hamiltonian \cite{PageWootters1983}. Resource theories of asymmetry organize the harmonic coherence required to break time-translation symmetry \cite{MarvianSpekkens2014Modes}, and shared asymmetry has been used to quantify resources in relational-time constructions \cite{CarmoSoaresPinto2021}. Quantum Fisher information likewise has an operational resource interpretation for energetic coherence \cite{Marvian2022,YamaguchiTajima2023}, while quantitative Wigner--Araki--Yanase and conservation-law results constrain operations that require a coherence reference \cite{MarvianSpekkensWAY2012,TajimaShiraishiSaito2020}. These frameworks are essential context, but they leave open a different local question. A clock and signal may exchange an exact energy gap while the total state remains stationary, so the global asymmetry can vanish identically even though a relative temporal mode is statistically resolvable.

A second obstruction appears if one tries to bound that mode directly by Fisher information and mean energy. We construct a two-level family whose mean excess energy is fixed while the Fisher information of an arbitrarily high Bohr-frequency tangent remains constant. What collapses is instead the radius of the affine parameter neighborhood over which that tangent remains a positive density operator. Thus a purely local metric does not retain enough state-space geometry to support an energy--frequency law. This observation leads to the first resource regime: when the affine physical tangent radius is nonzero, the robust combination of Fisher information and radius is bounded by the population of the spectral endpoints that support the transition. In an autonomous clock--signal exchange the same quantity is bounded on both local sides, even for a globally stationary baseline with pre-existing relational coherence.

The affine radius can, however, vanish at the boundary of state space while an exact nonlinear family remains physical. A first-order amplitude can enter a previously empty spectral sector provided the population needed for positivity is created at second order. Rank-changing quantum Fisher and Bures geometry is known to exhibit precisely such singular behavior \cite{Safranek2017}. Here we ask for its operational temporal-resource consequence. We show that the missing resource moves from zeroth-order spectral survival to positive second-order spectral synthesis. For one endpoint, arbitrary finite-copy collective measurements obey a Fisher--curvature bound. When opposite endpoint orientations participate in the same temporal mode, their score amplitudes combine by a sharp Minkowski law rather than by ordinary addition. In an exact autonomous exchange this produces a two-sided synthesis-action bound that remains nonzero inside a fixed-total-energy eigenspace, where both global time-translation asymmetry and signed total-energy curvature vanish.

The resulting picture is a two-regime resource principle for relative temporal information. At nonzero affine radius, information is backed by \emph{pre-existing spectral survival}; at a rank-changing zero-radius boundary, it is backed by \emph{spectral synthesis action}. Both statements hold for arbitrary finite copy number and arbitrary collective POVMs. Their coefficients are sharp in finite-dimensional fixed-total-energy models. We further show that the boundary law survives arbitrary coherent support through a support/endpoint geometric reduction and that multiple exchange frequencies share one positive Hessian budget; a single Fourier measurement simultaneously saturates the complete frequency-weighted sum in a mode-separated fixed-energy shell.

These results are complementary to, rather than replacements for, established asymmetry and metrology bounds. We do not propose a new generic resource theory of time, and the synthesis action is not identified with the total dynamical energy required to implement an encoding. Protocol-level energy-constrained metrology addresses that broader implementation cost \cite{ChenYang2026}, while waveform Holevo theory treats attainable information under general linear waveform estimation \cite{GardnerEtAl2024}. Our claim is narrower: for a finite autonomous system carrying an exact relative temporal mode, there are frequency-resolved necessary state-family resources that remain operative when global asymmetry is zero, and the relevant resource changes order at a rank-changing boundary.
